\documentclass[11pt]{article}

\usepackage[a4paper,margin=0.78in]{geometry}
\usepackage[T1]{fontenc}
\usepackage[utf8]{inputenc}
\usepackage{lmodern}
\usepackage{amsmath,amssymb,amsthm,mathtools}
\usepackage{enumitem,booktabs,array,tabularx,microtype,xcolor}
\usepackage[numbers,sort&compress]{natbib}
\usepackage[colorlinks=true,linkcolor=blue!65!black,
  citecolor=blue!65!black,urlcolor=blue!70!black]{hyperref}
\usepackage[nameinlink,noabbrev]{cleveref}
\setlength{\emergencystretch}{3em}
\setlist{itemsep=0pt,topsep=3pt,parsep=0pt,partopsep=0pt}

\hypersetup{
  pdftitle={Refining the H1 Crosswalk Frontier},
  pdfauthor={Liang Wang},
  pdfsubject={Three independent production roots after archive separation and formal gluing},
  pdfkeywords={fixed shift, occurrence crosswalk, blocker antichain, route decision}
}

\newtheorem{theorem}{Theorem}[section]
\newtheorem{proposition}[theorem]{Proposition}
\newtheorem{corollary}[theorem]{Corollary}
\theoremstyle{definition}
\newtheorem{definition}[theorem]{Definition}
\theoremstyle{remark}
\newtheorem{remark}[theorem]{Remark}

\newcommand{\Lzero}{\mathrm{L0}}
\newcommand{\Lone}{\mathrm{L1}}
\newcommand{\Ltwo}{\mathrm{L2}}
\newcommand{\NT}{\textnormal{\textsc{not-testable}}}
\newcommand{\PROVED}{\textnormal{\textsc{proved}}}
\newcommand{\Hmono}{\mathsf H_{\rm mono}}
\newcommand{\Hlocal}{\mathsf H_{\rm local}}
\newcommand{\Hsupport}{\mathsf H_{\rm support}}
\newcommand{\Hcanon}{\mathsf H_{\rm canon}}
\newcommand{\Hoverlap}{\mathsf H_{\rm overlap}}
\newcommand{\Htotal}{\mathsf H_{\rm total}}
\newcommand{\Hwitness}{\mathsf H_{\rm witness}}
\newcolumntype{Y}{>{\raggedright\arraybackslash}X}
\newcolumntype{P}[1]{>{\raggedright\arraybackslash}p{#1}}

\title{\textbf{Refining the H1 Crosswalk Frontier:}\\
\textbf{Three Independent Production Roots after}\\
\textbf{Archive Separation and Formal Gluing}}
\author{Liang Wang\\
School of Artificial Intelligence and Automation,\\
Huazhong University of Science and Technology, Wuhan 430074, P.R. China\\
\texttt{wangliang.f@gmail.com}}
\date{July 2026}

\begin{document}
\maketitle

\begin{abstract}
TPC-162 records the first missing structural object as the monolithic
node \(\Hmono\).
TPC-163--165 now reveal its internal dependency structure without
constructing the production carrier.  We freeze the TPC-162 pointer
as historical fact and refine it by a source-locked acyclic graph.

Inside this monolithic-crosswalk sub-DAG, the graph has exactly three
minimal \(\mathsf{NOT\_TESTABLE}\) roots:
\[
\begin{gathered}
\mathsf{H1.source\_backed\_local\_occurrence\_edge\_family},\\
\mathsf{H1.actual\_active\_support\_certificate},\\
\mathsf{H1.canonical\_minimal\_representation\_certificate}.
\end{gathered}
\]
The first root feeds overlap compatibility and formal
local-to-global totality.  Formal totality then joins the other two
roots at a production occurrence witness, which feeds the historical
monolithic crosswalk.  The exact archived separation key of TPC-164
and the finite gluing theorem of TPC-165 are proved supporting nodes,
not unresolved roots.

For work ordering, the local-edge family is selected as the refined
first-child pointer.  The full three-element antichain is nevertheless
preserved: selection is not implication.  TPC-163 still finds zero
theorem-backed production local occurrence edges, so the verdict
remains \(\mathsf{NOT\_TESTABLE}\).  No selected route is stopped, and
no actual-carrier impossibility is inferred.  These three roots are
not the complete H1 blocker set: the full map clause still has nine
zero-defect requirements and an independent occurrence registry,
while the scalar-plus-ETO clause remains independently
\(\mathsf{NOT\_TESTABLE}\).
\end{abstract}

\noindent\textbf{Keywords:} fixed shift; occurrence provenance;
dependency DAG; minimal blocker antichain; route decision.

\medskip
\noindent\fcolorbox{red!55!black}{red!3}{%
\begin{minipage}{0.96\linewidth}
\textbf{Claim boundary.}
This paper refines a frontier; it does not close it.  It does not
rewrite TPC-162, prove a production local occurrence family, certify
actual active support, prove a canonical/minimal actual
representation, complete the production crosswalk, stop the selected
route, establish positive fixed-\(X\) \(\Ltwo\), give a prime-pair
lower bound, or prove the twin-prime conjecture.  The three-root
antichain is scoped only to the monolithic-crosswalk sub-DAG, not to
the complete H1 contract.
\end{minipage}}

\section{Historical pointer versus current refinement}

TPC-162 freezes the current verdict
\[
 \mathsf{Verdict}_{162}=\NT
\]
and the first missing object
\[
 \Hmono:=
 \mathsf{H1.theorem\_backed\_occurrence\_provenance\_crosswalk}.
 \tag{1.1}\label{eq:historical}
\]
This statement is correct at the interface resolution used there
\citep{WangTPC162}.  A later decomposition should not retroactively
replace \eqref{eq:historical}: reproducible route histories require
old pointers to remain interpretable against their own source locks.

TPC-163--165 provide new resolution:
\begin{enumerate}[leftmargin=2.1em]
 \item the frozen source corpus has no theorem-backed production
       actual-occurrence edge, and native triples collide
       \citep{WangTPC163};
 \item a unique five-field minimum key separates every archived cut
       row \citep{WangTPC164}; and
 \item compatible local formal patches would glue uniquely with exact
       conservation, but no production patch is supplied
       \citep{WangTPC165}.
\end{enumerate}
These facts justify a new child-level graph beneath \(\Hmono\).

\subsection{The complete H1 envelope remains outside this refinement}

TPC-156 defines two sufficient clauses \citep{WangTPC156}.  Write
\[
\mathcal D=
\{D_L,D_{QD},D_{QZ},D_G,D_P,D_{DZ},D_{GP},
  D_{\rm cover},D_{\rm rec}\}.
\]
Then
\[
\begin{aligned}
\mathsf{Map}_{H1}
 &=
 \Hmono\wedge
 \bigwedge_{D\in\mathcal D}(D=0)
 \wedge\mathsf{OccurrenceRegistryTotal},\\
\mathsf{Scalar}_{H1}
 &=
 \mathsf{CompleteFUM}_{o(X)}
 \wedge\mathsf{TheoremBackedETO},\\
\mathsf{H1Complete}
 &=\mathsf{Map}_{H1}\vee\mathsf{Scalar}_{H1}.
\end{aligned}
\tag{1.2}\label{eq:full-h1}
\]
TPC-166 refines only the factor \(\Hmono\) in
\(\mathsf{Map}_{H1}\).  The nine zero-defect conditions and the
independent occurrence registry remain external requirements.  The
scalar-plus-ETO clause is a separate \(\NT\) route and is not a child
of \(\Hmono\).  Therefore the antichain computed below must never be
reported as the complete H1 minimal blocker antichain.

\begin{definition}[Non-retroactive refinement]\label{def:refine}
A refinement of a historical frontier node \(H\) consists of:
\begin{enumerate}[label=(\roman*),leftmargin=2.1em]
 \item a source lock preserving \(H\)'s original identifier, status,
       and scope;
 \item a directed acyclic graph whose target is a new copy of \(H\);
 \item typed child nodes whose conjunction supplies \(H\); and
 \item an explicit statement that the historical record is not
       rewritten.
\end{enumerate}
\end{definition}

\section{The refined production DAG}

We abbreviate:
\[
\begin{aligned}
\Hlocal&=\mathsf{source\mbox{-}backed\ local\ occurrence\ edge\ family},\\
\Hsupport&=\mathsf{actual\ active\ support\ certificate},\\
\Hcanon&=\mathsf{canonical/minimal\ representation\ certificate},\\
\Hoverlap&=\mathsf{local\ overlap\ bijection/cocycle},\\
\Htotal&=\mathsf{glued\ formal\ occurrence\ totality},\\
\Hwitness&=\mathsf{production\ occurrence\ witness}.
\end{aligned}
\]
The dependency direction is ``parent is prerequisite.''  The full
refined graph is:
\[
\begin{array}{ccl}
\Hlocal&\longrightarrow&\Hoverlap\\
\{\Hlocal,\Hoverlap,K_{164},G_{165}\}
 &\longrightarrow&\Htotal\\
\{\Htotal,\Hsupport,\Hcanon\}
 &\longrightarrow&\Hwitness\\
\Hwitness&\longrightarrow&\Hmono,
\end{array}
\tag{2.1}\label{eq:dag}
\]
where \(K_{164}\) is the proved archived separation key and
\(G_{165}\) is the proved finite typed gluing theorem.  The three
nodes \(\Hlocal,\Hsupport,\Hcanon\) have no parents.  The nodes
\(K_{164},G_{165}\) are also parentless, but their status is
\(\PROVED\), so they are not blockers.

\begin{definition}[Minimal unresolved root antichain]\label{def:antichain}
For a target \(H\) in a finite prerequisite DAG, take the ancestral
subgraph of \(H\).  An unresolved root is a node of status \(\NT\)
having no \(\NT\) parent in that subgraph.  The set of all such roots
is the \emph{minimal unresolved root antichain}.
\end{definition}

This definition removes descendants that are unavailable only because
an earlier unresolved prerequisite is missing.  It does not choose
one root at the expense of independent roots.

\section{Exact three-root frontier}

\begin{theorem}[Refined crosswalk-sub-DAG blocker antichain]
\label{thm:antichain}
Within the source-locked DAG \eqref{eq:dag}, the minimal unresolved
root antichain of \(\Hmono\) is exactly
\[
 \mathcal A_{166}
 =
 \{\Hlocal,\Hsupport,\Hcanon\}.
 \tag{3.1}\label{eq:antichain}
\]
The nodes \(\Hoverlap,\Htotal,\Hwitness,\Hmono\) are unresolved
descendants and hence do not belong to \(\mathcal A_{166}\).
\end{theorem}

\begin{proof}
The imported TPC-164 and TPC-165 audits give
\[
 \operatorname{status}(K_{164})
 =\operatorname{status}(G_{165})=\PROVED.
\]
TPC-163 gives zero qualifying production local edges, and TPC-165
keeps formal totality, actual support, and canonical/minimality
separate.  Therefore
\[
 \operatorname{status}(\Hlocal)
 =\operatorname{status}(\Hsupport)
 =\operatorname{status}(\Hcanon)=\NT.
\]
Each of these three has no unresolved parent, so all belong to the
root antichain.

The node \(\Hoverlap\) has unresolved parent \(\Hlocal\).
The node \(\Htotal\) has unresolved parents \(\Hlocal,\Hoverlap\).
The node \(\Hwitness\) has unresolved parents
\(\Htotal,\Hsupport,\Hcanon\), and \(\Hmono\) has unresolved parent
\(\Hwitness\).  Hence none of these descendants is minimal.  The DAG
contains no other nodes, proving \eqref{eq:antichain}.
\end{proof}

\begin{remark}[Not the full H1 antichain]\label{rem:not-full}
\Cref{thm:antichain} concerns only the ancestral sub-DAG of
\(\Hmono\).  Replacing \(\Hmono\) by its three roots inside
\eqref{eq:full-h1} still leaves the nine defect nodes,
occurrence-registry totality, and the independent scalar-plus-ETO
clause.  Their full-DAG minimality is not computed here.
\end{remark}

\begin{corollary}[No single-root reduction]\label{cor:no-single}
Proving any one member of \(\mathcal A_{166}\) does not, by the
current dependency contract, prove either of the other two.
\end{corollary}

\begin{proof}
The three nodes are pairwise parent-independent in
\eqref{eq:dag}.  No imported theorem supplies a directed path between
them.
\end{proof}

\section{A selected pointer that preserves the antichain}

For sequential work, a single display pointer is useful.  We select
\[
 \mathsf{FirstChild}_{166}=\Hlocal
 \tag{4.1}\label{eq:selected}
\]
because it is the first source-producing object on the
local-to-global branch.  If local edges are found, overlap maps become
testable and the proved gluing theorem can be instantiated.

\begin{proposition}[Selection is not erasure]\label{prop:selection}
Equation \eqref{eq:selected} is an ordering convention for the next
construction.  The production blocker set remains the full
\(\mathcal A_{166}\) of \cref{eq:antichain}.
\end{proposition}

\begin{proof}
The machine decision stores both fields independently:
\[
\begin{split}
\mathsf{canonical\ selected\ representative}&=\Hlocal,\\
\mathsf{full\ antichain\ preserved}
 &= [\Hlocal,\Hsupport,\Hcanon].
\end{split}
\]
Its validator recomputes the antichain from the DAG and rejects any
candidate that deletes \(\Hsupport\) or \(\Hcanon\).  Thus selection
cannot change the mathematical dependency set.
\end{proof}

Here ``canonical selected representative'' refers only to deterministic
serialization and display order.  It is not the missing
\(\Hcanon\) theorem.

\section{Production evidence and verdict}

The source-locator census remains decisive at the first child:
\[
 \#\mathcal E_{\rm local}^{\rm theorem\mbox{-}backed}=0
 \quad\text{in the frozen declared corpus}.
 \tag{5.1}\label{eq:zero-local}
\]
Consequently no production local patch family exists in the current
evidence set.  The statuses are:
\begin{center}
\begin{tabularx}{0.96\linewidth}{@{}Y P{0.2\linewidth}@{}}
\toprule
\textbf{Object} & \textbf{Status}\\
\midrule
minimum archived separation key & \(\PROVED\)\\
finite typed gluing theorem & \(\PROVED\)\\
source-backed local occurrence edge family & \(\NT\)\\
actual active-support certificate & \(\NT\)\\
canonical/minimal representation certificate & \(\NT\)\\
overlap cocycle and glued formal totality & \(\NT\)\\
production occurrence witness and monolithic crosswalk & \(\NT\)\\
\bottomrule
\end{tabularx}
\end{center}

\begin{theorem}[Current refined decision]\label{thm:decision}
Under the nine source locks of TPC-166, and only inside the historical
monolithic-crosswalk sub-DAG,
\[
 \boxed{\mathsf{Verdict}_{166}=\NT},
 \qquad
 \mathsf{FirstChild}_{166}=\Hlocal,
 \qquad
 \mathcal A_{166}
 =\{\Hlocal,\Hsupport,\Hcanon\}.
 \tag{5.2}\label{eq:decision}
\]
The selected occurrence-augmented route is not stopped, and no
actual-carrier impossibility is proved.
\end{theorem}

\begin{proof}
The verdict cannot be \(\mathsf{GO}\) because each root in
\(\mathcal A_{166}\) is unresolved.  It cannot be a route stop or
impossibility verdict because the source census is closed-world and
the formal gluing theorem leaves source augmentation open.  Thus the
typed result is \(\NT\), with pointer and antichain given by
\cref{eq:selected,eq:antichain}.
\end{proof}

\section{Why TPC-162 is not rewritten}

There is no contradiction between
\[
 \mathsf{FirstMissing}_{162}=\Hmono
 \quad\text{and}\quad
 \mathsf{FirstChild}_{166}=\Hlocal.
\]
The first names the unresolved public interface at TPC-162
resolution.  The second names one child after a later factorization.
The certificate records
\[
 \mathsf{historical\ pointer\ retroactively\ rewritten}=\mathsf{false}
\]
and mutation-tests the opposite value.  This provenance discipline
allows later papers to refine a route without silently changing what
an earlier source actually asserted.

\section{Reproducible DAG audit}

The generator locks the TPC-162 snapshot/audit and both
certificate/audit pairs from TPC-163--165, together with the TPC-156
full H1 contract.  It validates:
\begin{enumerate}[leftmargin=2.1em]
 \item unique node identifiers and total parent references;
 \item acyclicity by depth-first topological traversal;
 \item exact node statuses and dependency lists;
 \item recomputation of the minimal unresolved root antichain;
 \item preservation of the crosswalk-sub-DAG scope and exclusion of
       the nine defects, independent registry, and scalar clause;
 \item preservation of the historical pointer;
 \item zero production local edges and the \(\NT\) verdict.
\end{enumerate}
Real mutations delete either independent root, rewrite the historical
pointer, fabricate a production edge, promote the verdict to
\(\mathsf{GO}\), stop the route, introduce a dependency cycle, or
select a representative outside the antichain.  The same validator
rejects each mutation.

\section{Next source-producing task}

The refinement identifies what the next paper must actually deliver:
not another global schema, but at least one theorem-backed local
actual-occurrence edge family with:
\begin{enumerate}[leftmargin=2.1em]
 \item the five-field archived source address;
 \item canonical path and hash, theorem and formula locators, and a
       nonempty derivation tree;
 \item exact typed row weights and parent/stage/downstream payloads;
 \item an explicit statement of whether the result is merely formal
       or lies on the actual carrier.
\end{enumerate}
A partial local theorem would be genuine forward progress even before
total coverage.  It would not close \(\Hsupport\) or \(\Hcanon\);
those remain parallel research programs.

\section{Conclusion}

TPC-166 turns one opaque H1 blocker into a precise three-root frontier
inside that blocker.  Two supporting uncertainties are already closed:
archived
cuts have an exact lossless key, and compatible formal patches glue
with exact conservation.  The unresolved mathematics is now
localized to source-backed local occurrence edges, actual active
support, and canonical/minimal representation.  The route remains
open, the verdict remains \(\NT\), and the next constructive target
is smaller without being overstated.  The complete H1 graph,
including the external map-clause and scalar-clause requirements of
\eqref{eq:full-h1}, remains for a later integration.

\bibliographystyle{plainnat}
{\small\bibliography{references}}

\end{document}
